\textbf{COLUNI 2021 -- Questão 30}\par\medskip Leia o texto.

\par

[...] As bactérias coçavam agradavelmente a pele da terra ao maquiá-la com suas geleias, para que ficasse novamente boa e porosa. A palha de milho cobria ainda em camada grossa o chão e protegia a terra contra o calor do sol e a força impetuosa das chuvas e garantia umidade suficiente para toda essa vida ativa que rodopiava e formigava nos minitúneis dos poros. Até pequenas minhocas apareciam fazendo suas galerias, brincando com os nematoides dos quais elas não gostavam. Mas tudo isso ocorria somente na camada superficial. Em pouca profundidade, a vida acabava. Aqui reinava um silêncio intenso e absoluto que a laje dura, a falta de ar e a presença de dezenas de antibióticos impunham. Estes foram lavados da camada superficial e banidos para essas regiões. Os fungos que os produziam queriam, com isso, assegurar a sua supremacia na terra e não queriam entender que seus antibióticos não podiam ficar na camada superficial densamente povoada. [...]

\par

PRIMAVESI, A. A terra e o arado. IN: PRIMAVESI, A. \textbf{A convenção dos ventos e outros contos atuais do mundo em que vivemos}. 2. ed. São Paulo: Expressão Popular, 2016. p. 89.

\par

O solo apresenta grande diversidade de organismos que interagem e realizam diversas funções contribuindo para saúde desse ambiente. Sobre a diversidade apresentada no texto, é INCORRETO afirmar que:\begin{enumerate}
\item As bactérias são organismos unicelulares e procariontes, importantes decompositores da matéria orgânica
       no solo e podem estabelecer relações íntimas com as plantas, promovendo a fixação de nitrogênio.
\item As minhocas, animais pertencentes ao Filo Annelida, são hermafroditas com circulação fechada e
       desempenham importante papel ao humificar o solo, contribuindo para a nutrição das plantas.
\item Os nematoides, vermes de corpo cilíndrico, pertencentes ao Filo Nematoda, apresentam sistema digestório
       incompleto e podem causar danos à agricultura por parasitar raízes de plantas.
\item Os fungos são organismos eucariotos heterótrofos, importantes decompositores, que podem se associar a
       raízes de plantas formando as micorrizas.

\end{enumerate}
